\documentclass[12pt,a4paper]{article}

\usepackage[utf8]{inputenc}
\usepackage[T1]{fontenc}
\usepackage[romanian]{babel}
\usepackage[a4paper, margin=2.5cm]{geometry}
\usepackage{graphicx}
\usepackage{amsmath, amssymb}
\usepackage{booktabs}
\usepackage{hyperref}
\usepackage{xcolor}
\usepackage{listings}
\usepackage{float}
\usepackage{caption}
\usepackage{array}
\usepackage{enumitem}

\hypersetup{
    colorlinks=true,
    linkcolor=black,
    urlcolor=blue,
    citecolor=blue
}

\lstdefinestyle{pystyle}{
    language=Python,
    basicstyle=\ttfamily\footnotesize,
    keywordstyle=\color{blue}\bfseries,
    commentstyle=\color{gray}\itshape,
    stringstyle=\color{red!70!black},
    numbers=left,
    numberstyle=\tiny\color{gray},
    breaklines=true,
    frame=single,
    backgroundcolor=\color{gray!5},
    showstringspaces=false,
    tabsize=2
}
\lstset{style=pystyle}

\title{
    \vspace{2cm}
    \textbf{Recomandarea produselor: abordări bazate pe \\ Machine Learning} \\
    \vspace{0.5cm}
    \large Sistem de recomandare pentru filme bazat pe \\ algoritmi de învățare automată
}
\author{Vînt Alexandru}
\date{Mai 2026}

\begin{document}

\maketitle
\thispagestyle{empty}
\vspace{2cm}

\begin{center}
\textit{Disciplina:} Sisteme Inteligente \\
\textit{Set de date:} MovieLens (ml-latest-small)
\end{center}

\newpage
\tableofcontents
\newpage

% =====================================================================
\section{Introducere}

Sistemele de recomandare au devenit o componentă esențială a platformelor digitale moderne (Netflix, YouTube, Amazon, Spotify), influențând semnificativ atât experiența utilizatorilor, cât și performanța economică a serviciilor respective. Cantitatea masivă de conținut disponibil face imposibilă explorarea manuală a tuturor opțiunilor, motiv pentru care algoritmii de recomandare automatizează procesul de filtrare și sugerare a elementelor relevante pentru fiecare utilizator.

\subsection{Obiectivul proiectului}

Obiectivul acestui proiect este dezvoltarea unui sistem de recomandare pentru filme care sugerează utilizatorilor titluri în funcție de preferințele și istoricul lor de vizionare, folosind tehnici diverse de \textit{Machine Learning}. Sunt implementate, evaluate și comparate șapte abordări distincte, acoperind atât metode clasice (filtrare colaborativă bazată pe vecini, modele bazate pe conținut), cât și metode moderne (matrix factorization, deep learning).

\subsection{Tehnologii utilizate}

\begin{itemize}[noitemsep]
    \item \textbf{Limbaj}: Python 3
    \item \textbf{Manipulare date}: pandas, NumPy
    \item \textbf{Modele clasice ML}: scikit-learn (k-NN, Random Forest, TF-IDF, metrici)
    \item \textbf{Gradient boosting}: XGBoost
    \item \textbf{Deep Learning}: PyTorch
    \item \textbf{Vizualizare}: Matplotlib, Seaborn
    \item \textbf{Mediu de dezvoltare}: Jupyter Notebook
\end{itemize}

\subsection{Structura proiectului}

\begin{itemize}[noitemsep]
    \item \texttt{src/python.ipynb} -- notebook-ul principal cu toate experimentele
    \item \texttt{src/dataset/} -- setul de date MovieLens (ratings.csv, movies.csv, tags.csv, links.csv)
    \item \texttt{documentatie/} -- raportul de proiect
    \item \texttt{screenshots/} -- capturi de ecran cu rezultate intermediare
\end{itemize}

% =====================================================================
\section{Setul de date}

\subsection{Sursa datelor}

Datele provin din proiectul \textbf{MovieLens} dezvoltat de grupul de cercetare GroupLens din cadrul Departamentului de Informatică al Universității din Minnesota. A fost folosită versiunea \textit{ml-latest-small}, descărcată în data de 17 martie 2026 de la adresa \url{https://grouplens.org/datasets/movielens/latest/}.

\subsection{Conținutul setului de date}

Setul de date conține:
\begin{itemize}[noitemsep]
    \item \textbf{100.836} de rating-uri
    \item \textbf{3.683} de tag-uri textuale
    \item \textbf{9.742} de filme
    \item \textbf{610} utilizatori (fiecare cu cel puțin 20 de rating-uri)
\end{itemize}

\subsection{Fișierele utilizate}

\begin{table}[H]
\centering
\begin{tabular}{lll}
\toprule
\textbf{Fișier} & \textbf{Coloane} & \textbf{Descriere} \\
\midrule
\texttt{ratings.csv} & userId, movieId, rating, timestamp & Rating-uri 0.5--5.0 stele \\
\texttt{movies.csv} & movieId, title, genres & Titlu + genuri (pipe-separated) \\
\texttt{tags.csv} & userId, movieId, tag, timestamp & Tag-uri text aplicate de utilizatori \\
\texttt{links.csv} & movieId, imdbId, tmdbId & Identificatori externi \\
\bottomrule
\end{tabular}
\caption{Structura setului de date MovieLens}
\end{table}

% =====================================================================
\section{Preprocesarea datelor}

\subsection{Tratarea valorilor lipsă pentru genuri}

Filmele care nu au genuri asociate sunt marcate cu valoarea specială \texttt{(no genres listed)}. Aceasta a fost înlocuită cu un șir vid și ulterior s-a creat o coloană \texttt{genres\_list} obținută prin separarea genurilor pe caracterul \texttt{|}.

\begin{lstlisting}
movies['genres'] = movies['genres'].replace('(no genres listed)', '')
movies['genres_list'] = movies['genres'].apply(
    lambda x: x.split('|') if x != '' else []
)
\end{lstlisting}

\subsection{Extragerea anului de lansare}

Anul de lansare al fiecărui film este prezent în titlu, între paranteze. A fost extras cu o expresie regulată și mutat într-o coloană separată \texttt{year}:

\begin{lstlisting}
movies['year'] = movies['title'].str.extract(r'\((\d{4})\)').astype(float)
movies['title'] = movies['title'].str.replace(r'\s*\(\d{4}\)', '', regex=True)
\end{lstlisting}

Pentru 13 filme la care anul lipsea, valorile au fost completate manual pe baza titlului (\texttt{year\_map}).

\subsection{Conversia timestamp-urilor și normalizarea tag-urilor}

Coloanele \texttt{timestamp} din \texttt{ratings} și \texttt{tags} au fost convertite din \textit{Unix epoch} în obiecte \texttt{datetime} pentru a permite extragerea de features temporale. Tag-urile au fost normalizate la minuscule, eliminându-se spațiile inutile.

% =====================================================================
\section{Analiza exploratorie a datelor (EDA)}

\subsection{Distribuția rating-urilor}

Rating-urile urmează o distribuție ușor înclinată spre stânga (\textit{left-skewed}): utilizatorii tind să acorde note mai mari decât cele mici. Valoarea cea mai frecventă este \textbf{4.0}, urmată de \textbf{3.0} și \textbf{5.0}. Rating-urile cu jumătăți de stea (0.5, 1.5, 2.5, etc.) sunt semnificativ mai rare.

\subsection{Statistici pe utilizatori și filme}

\begin{itemize}[noitemsep]
    \item Numărul de rating-uri pe utilizator are o distribuție puternic \textit{long-tail}: cei mai activi utilizatori au câteva mii de rating-uri, în timp ce mediana este mult mai mică.
    \item La fel, distribuția numărului de rating-uri pe film este puternic dezechilibrată: câteva filme populare acumulează majoritatea voturilor, iar marea masă a filmelor are puține rating-uri.
\end{itemize}

\subsection{Matricea de corelație}

Analiza corelațiilor între variabilele numerice (rating, year, avg\_rating, num\_ratings, user\_avg\_rating, user\_num\_ratings) arată că cei mai puternici predictori pentru \texttt{rating} sunt:
\begin{itemize}[noitemsep]
    \item \texttt{avg\_rating} -- media rating-urilor primite de film
    \item \texttt{user\_avg\_rating} -- media rating-urilor date de utilizator
\end{itemize}

Aceste două variabile vor sta la baza modelelor supervizate (Random Forest, XGBoost).

% =====================================================================
\section{Feature Engineering}

Pe lângă variabilele brute, au fost construite mai multe features derivate care îmbogățesc reprezentarea perechilor (user, film):

\subsection{Statistici agregate}

\begin{itemize}[noitemsep]
    \item \textbf{Per film}: \texttt{avg\_rating}, \texttt{num\_ratings}
    \item \textbf{Per utilizator}: \texttt{user\_avg\_rating}, \texttt{user\_num\_ratings}
\end{itemize}

\subsection{One-hot encoding pentru genuri}

Pentru cele șapte genuri cele mai populare (Action, Comedy, Drama, Horror, Romance, Sci-Fi, Thriller) s-au creat coloane binare \texttt{is\_Action}, \texttt{is\_Comedy}, etc.

\subsection{Features temporale}

Din timestamp-ul rating-ului au fost extrase:
\begin{itemize}[noitemsep]
    \item \texttt{rating\_year}, \texttt{rating\_month}, \texttt{rating\_dayofweek}
    \item \texttt{years\_since\_release} = rating\_year $-$ year
\end{itemize}

\subsection{User-genre affinity}

Cea mai informativă caracteristică derivată: media rating-urilor pe care un utilizator le-a acordat filmelor dintr-un anumit gen. Operația constă în explodarea coloanei \texttt{genres\_list}, urmată de un pivot \texttt{userId} $\times$ \texttt{genre} cu agregare \texttt{mean}. Pentru combinațiile lipsă (utilizatori care nu au votat un gen) s-a folosit media generală a utilizatorului ca valoare default.

% =====================================================================
\section{Împărțirea datelor train/test}

Setul de rating-uri a fost împărțit aleator în \textbf{80\% antrenament} și \textbf{20\% testare}, cu \texttt{random\_state=42} pentru reproductibilitate.

\begin{lstlisting}
train_ratings, test_ratings = train_test_split(
    ratings_cf, test_size=0.2, random_state=42
)
\end{lstlisting}

Pentru evaluarea modelelor s-au folosit două metrici standard:
\begin{itemize}[noitemsep]
    \item \textbf{RMSE} (Root Mean Squared Error):
    $$\text{RMSE} = \sqrt{\frac{1}{N}\sum_{i=1}^{N}(\hat{r}_i - r_i)^2}$$
    \item \textbf{MAE} (Mean Absolute Error):
    $$\text{MAE} = \frac{1}{N}\sum_{i=1}^{N}|\hat{r}_i - r_i|$$
\end{itemize}

% =====================================================================
\section{Modele de recomandare}

În continuare sunt prezentate cele șapte abordări implementate.

% ---------------------------------------------------------------------
\subsection{k-NN Item-Based (filtrare colaborativă)}

\textbf{Ideea}: două filme sunt similare dacă au primit rating-uri asemănătoare de la aceiași utilizatori. Pentru a recomanda filme similare unui film $m$, se găsesc cei mai apropiați $k$ vecini ai săi într-un spațiu vectorial.

\textbf{Implementare}:
\begin{enumerate}[noitemsep]
    \item S-a construit matricea sparse \texttt{movieId $\times$ userId} cu valori rating.
    \item S-a antrenat \texttt{NearestNeighbors} cu metrica \textbf{cosine} și algoritm \texttt{brute}.
    \item Pentru un film dat, se obțin top-$k$ vecini, excluzând filmul însuși.
\end{enumerate}

\textbf{Recomandarea} returnează lista de filme cu cele mai mici distanțe cosinus.

% ---------------------------------------------------------------------
\subsection{k-NN User-Based}

\textbf{Ideea}: doi utilizatori sunt similari dacă au evaluat asemănător filmele comune. Recomandările pentru utilizatorul $u$ se obțin agregând rating-urile celor mai similari $k$ vecini, ponderat cu similaritatea.

\textbf{Algoritm}:
\begin{enumerate}[noitemsep]
    \item Construire matrice \texttt{userId $\times$ movieId}.
    \item Calcul vecini cu cosine similarity.
    \item Pentru fiecare film nevizionat de $u$:
    $$\hat{r}_{u,m} = \frac{\sum_{v \in N(u)} \text{sim}(u,v) \cdot r_{v,m}}{\sum_{v \in N(u)} \text{sim}(u,v)}$$
    \item Sortare descrescătoare după scor, returnare top-$N$.
\end{enumerate}

% ---------------------------------------------------------------------
\subsection{k-NN Content-Based cu TF-IDF}

\textbf{Ideea}: două filme sunt similare dacă tag-urile lor textuale sunt similare. Spre deosebire de abordările colaborative, această metodă \textbf{nu folosește rating-urile}, ci doar conținutul (tag-urile aplicate de utilizatori).

\textbf{Implementare}:
\begin{enumerate}[noitemsep]
    \item Agregarea tag-urilor per film: fiecare film devine un ,,document'' textual.
    \item Vectorizare TF-IDF cu \texttt{min\_df=2} (ignoră tag-uri foarte rare) și \texttt{max\_df=0.8} (ignoră tag-uri foarte comune).
    \item Antrenare k-NN cu metrica cosine pe matricea TF-IDF.
\end{enumerate}

\textbf{Formula TF-IDF} (pentru tag $t$ în film $m$):
$$\text{tfidf}(t, m) = \text{tf}(t, m) \cdot \log\frac{N}{\text{df}(t)}$$

Această metodă oferă explicabilitate excelentă: pentru fiecare recomandare se pot afișa tag-urile comune cu filmul sursă.

% ---------------------------------------------------------------------
\subsection{Random Forest Regressor}

\textbf{Ideea}: problema este transformată din \textit{recomandare} în \textit{regresie supervizată}. Pentru fiecare pereche (user, film) din setul de antrenament se prezice rating-ul folosind un ansamblu de arbori de decizie.

\textbf{Configurație finală (v3)}:
\begin{itemize}[noitemsep]
    \item Features: \texttt{avg\_rating}, \texttt{num\_ratings}, \texttt{user\_avg\_rating}, \texttt{user\_num\_ratings}, \texttt{year}, 7 features binare pentru genuri, 3 features temporale, 7 features de tip \texttt{affinity\_<gen>}.
    \item Hiperparametri: \texttt{n\_estimators=100}, \texttt{max\_depth=10}.
\end{itemize}

\textbf{Recomandarea} se generează evaluând modelul pe toate filmele nevizionate de un utilizator și sortând descrescător după rating-ul prezis. S-a aplicat un filtru suplimentar \texttt{num\_ratings >= 20} pentru a evita filmele obscure cu rating-uri inflaționiste.

% ---------------------------------------------------------------------
\subsection{XGBoost Regressor}

XGBoost este o implementare optimizată a algoritmului \textit{gradient boosting}, în care fiecare arbore corectează erorile arborilor anteriori. Comparativ cu Random Forest:
\begin{itemize}[noitemsep]
    \item Arborii sunt construiți \textbf{secvențial}, nu independent.
    \item Optimizarea folosește gradient descent pe o funcție obiectiv diferențiabilă.
    \item Regularizarea L1/L2 este integrată în funcția obiectiv.
\end{itemize}

\textbf{Hiperparametri folosiți}: \texttt{n\_estimators=150}, \texttt{max\_depth=6}, \texttt{learning\_rate=0.1}.

\textbf{Features}: \texttt{features\_extended} (fără \texttt{affinity\_<gen>}).

% ---------------------------------------------------------------------
\subsection{Matrix Factorization (SGD)}

\textbf{Ideea}: matricea rating-urilor $R \in \mathbb{R}^{|U|\times|M|}$ este aproximată prin produsul a două matrice cu rang mic:
$$R \approx P Q^T + \mu + b_u + b_i$$
unde:
\begin{itemize}[noitemsep]
    \item $P \in \mathbb{R}^{|U|\times k}$ -- matricea factorilor latenți pentru utilizatori
    \item $Q \in \mathbb{R}^{|M|\times k}$ -- matricea factorilor latenți pentru filme
    \item $\mu$ -- media globală
    \item $b_u$, $b_i$ -- bias-uri pentru utilizatori și filme
\end{itemize}

\textbf{Antrenare prin SGD}: parametrii se actualizează iterativ minimizând funcția de cost regularizată:
$$\mathcal{L} = \sum_{(u,i,r) \in \mathcal{T}} (r_{ui} - \hat{r}_{ui})^2 + \lambda \left(\|P_u\|^2 + \|Q_i\|^2 + b_u^2 + b_i^2\right)$$

\textbf{Regulile de actualizare} (pentru fiecare exemplu de antrenament):
\begin{align*}
e_{ui} &= r_{ui} - \hat{r}_{ui} \\
P_u &\leftarrow P_u + \eta (e_{ui} \cdot Q_i - \lambda P_u) \\
Q_i &\leftarrow Q_i + \eta (e_{ui} \cdot P_u - \lambda Q_i) \\
b_u &\leftarrow b_u + \eta (e_{ui} - \lambda b_u) \\
b_i &\leftarrow b_i + \eta (e_{ui} - \lambda b_i)
\end{align*}

\textbf{Hiperparametri folosiți}: \texttt{n\_factors=50}, \texttt{lr=0.005}, \texttt{reg=0.02}, \texttt{n\_epochs=15}.

\textbf{Notă importantă}: implementarea inițială folosea SVD trunchiat din \texttt{scipy.sparse.linalg.svds}, dar a fost înlocuită cu o implementare SGD personalizată, mult mai stabilă pe matrici extrem de sparse (RMSE: 3.16 $\rightarrow$ 0.88).

% ---------------------------------------------------------------------
\subsection{Deep Learning: Neural Collaborative Filtering (NCF)}

\textbf{Ideea}: în loc să folosim produs scalar liniar între factorii latenți (ca în MF), folosim o rețea neuronală pentru a învăța o funcție de interacțiune non-liniară între user și film.

\textbf{Arhitectură}:
\begin{enumerate}[noitemsep]
    \item \textbf{Embedding-uri}: \texttt{user\_embedding} și \texttt{item\_embedding}, fiecare de dimensiune 50.
    \item \textbf{Concatenare}: vectorul de input este \texttt{[user\_emb; item\_emb]} de dimensiune 100.
    \item \textbf{MLP}: trei straturi fully-connected (100 $\rightarrow$ 64 $\rightarrow$ 32 $\rightarrow$ 16), cu activare ReLU și dropout (p=0.2).
    \item \textbf{Strat de ieșire}: liniar (16 $\rightarrow$ 1).
\end{enumerate}

\textbf{Configurație antrenare}:
\begin{itemize}[noitemsep]
    \item Optimizator: Adam, \texttt{lr=0.005}, \texttt{weight\_decay=1e-5} (L2)
    \item Funcție de cost: MSE
    \item Batch size: 256
    \item Epoci: 10
\end{itemize}

\textbf{Evoluția loss-ului pe parcursul antrenării}:
\begin{table}[H]
\centering
\begin{tabular}{cc}
\toprule
\textbf{Epocă} & \textbf{Loss MSE} \\
\midrule
1 & 1.5035 \\
3 & 0.7865 \\
5 & 0.6898 \\
7 & 0.6306 \\
10 & 0.5382 \\
\bottomrule
\end{tabular}
\caption{Evoluția erorii pe parcursul antrenării NCF}
\end{table}

% =====================================================================
\section{Rezultate și comparație}

Toate modelele bazate pe regresie au fost evaluate pe același split train/test (80/20, \texttt{random\_state=42}).

\begin{table}[H]
\centering
\begin{tabular}{lcc}
\toprule
\textbf{Model} & \textbf{RMSE} & \textbf{MAE} \\
\midrule
XGBoost                          & \textbf{0.7790} & \textbf{0.5870} \\
Random Forest (v3, cu affinity)  & 0.7829 & 0.5905 \\
Matrix Factorization (SGD)       & 0.8820 & 0.6748 \\
Deep Learning (NCF)              & 0.9071 & 0.7124 \\
\bottomrule
\end{tabular}
\caption{Comparația performanței modelelor de regresie}
\label{tab:rezultate}
\end{table}

\subsection{Discuție}

\textbf{XGBoost} obține cel mai bun scor, fiind urmat foarte aproape de Random Forest v3. Diferența mică (0.004 RMSE) între cele două sugerează că, pentru acest set de date, contribuția suplimentară a boosting-ului față de bagging este modestă. Ambele modele beneficiază din plin de \textbf{feature engineering-ul} manual (în special \texttt{avg\_rating} și \texttt{user\_avg\_rating}, care domină feature importance-ul).

\textbf{Matrix Factorization} oferă un rezultat decent (RMSE 0.88) folosind exclusiv informația din matricea de rating-uri, fără features derivate. Acest lucru evidențiază puterea metodelor pure de filtrare colaborativă pe seturi de date mari și dense.

\textbf{NCF} obține cel mai slab scor în acest experiment. Posibile cauze:
\begin{itemize}[noitemsep]
    \item Setul de date este relativ mic (100k rating-uri) pentru a antrena o rețea cu sute de mii de parametri.
    \item Doar 10 epoci pot fi insuficiente; loss-ul de antrenament continuă să scadă.
    \item Lipsa unei strategii de \textit{early stopping} sau validare pe un subset separat.
\end{itemize}

\textbf{Modelele k-NN} (item-based, user-based, content-based) generează direct recomandări fără un scor de rating prezis, deci nu apar în tabelul comparativ cu RMSE/MAE. Ele oferă în schimb explicabilitate (vecinii pot fi inspectați) și sunt utile pentru \textbf{cold-start} (în special content-based, care nu necesită rating-uri).

% =====================================================================
\section{Concluzii}

Proiectul a explorat o gamă largă de tehnici pentru construirea unui sistem de recomandare de filme, de la metode clasice (k-NN, TF-IDF) până la abordări moderne (Matrix Factorization, Neural Collaborative Filtering). Principalele observații sunt:

\begin{enumerate}[noitemsep]
    \item \textbf{Feature engineering-ul} contează semnificativ: cele mai performante modele (XGBoost, RF) folosesc features derivate manual (\texttt{avg\_rating}, \texttt{user\_avg\_rating}, \texttt{user-genre affinity}).
    \item \textbf{Modelele simple performează bine}: pe un set de date de dimensiuni modeste, modelele supervizate clasice depășesc abordările bazate pe deep learning.
    \item \textbf{Alegerea metricii potrivite} este crucială: RMSE și MAE măsoară precizia predicției rating-urilor, dar nu reflectă neapărat calitatea ordonării recomandărilor (Top-N) percepute de utilizator.
    \item \textbf{Diversitatea abordărilor} are valoare practică: într-un sistem real, k-NN content-based ar rezolva problema cold-start (filme noi fără rating-uri), MF/NCF ar acoperi cazul general, iar XGBoost ar oferi predicții fine de rating.
\end{enumerate}

\subsection{Direcții viitoare}

\begin{itemize}[noitemsep]
    \item Evaluare pe metrici Top-N (Precision@K, Recall@K, NDCG).
    \item Implementare \textit{hybrid recommender} care combină predicțiile mai multor modele.
    \item Tuning extins de hiperparametri (Optuna / GridSearchCV).
    \item Strategii de \textit{early stopping} și \textit{learning rate scheduling} pentru NCF.
    \item Integrarea explicită a tag-urilor TF-IDF în features-urile XGBoost.
\end{itemize}

% =====================================================================
\section*{Bibliografie}
\addcontentsline{toc}{section}{Bibliografie}

\begin{enumerate}[noitemsep]
    \item F. Maxwell Harper și Joseph A. Konstan, ,,The MovieLens Datasets: History and Context'', \textit{ACM Transactions on Interactive Intelligent Systems}, vol. 5, nr. 4, 2015. \url{https://doi.org/10.1145/2827872}
    \item Y. Koren, R. Bell, C. Volinsky, ,,Matrix Factorization Techniques for Recommender Systems'', \textit{IEEE Computer}, 2009.
    \item X. He et al., ,,Neural Collaborative Filtering'', \textit{Proceedings of WWW 2017}.
    \item T. Chen, C. Guestrin, ,,XGBoost: A Scalable Tree Boosting System'', \textit{KDD 2016}.
    \item L. Breiman, ,,Random Forests'', \textit{Machine Learning}, vol. 45, 2001.
    \item Documentația scikit-learn: \url{https://scikit-learn.org/stable/}
    \item Documentația PyTorch: \url{https://pytorch.org/docs/}
\end{enumerate}

\end{document}
